\documentclass[a4paper, oneside]{article}
\usepackage{graphicx}
\usepackage{amsthm}
\usepackage{amsmath}
\usepackage{amssymb}
\usepackage[a4paper,
            bindingoffset=0.2in,
            left=2cm,
            right=2cm,
            top=2cm,
            bottom=2cm,
            footskip=.25in]{geometry}
\usepackage[italian]{babel}
\usepackage{pgfplots}
\usepackage{tabularx}
\usepackage{tikz}
\usepackage{wrapfig}
\usepackage{color}
\usepackage[d]{esvect}
\usepackage{circuitikz}
\definecolor{page}{rgb}{0.129,0.157,0.212}
\pagecolor{page}
\color{white}
\graphicspath{ {./images/} }
\usetikzlibrary{shapes.geometric}
\usetikzlibrary{datavisualization}
\usetikzlibrary{datavisualization.formats.functions}
\usetikzlibrary{patterns}
\pgfplotsset{width=10cm,compat=1.18}

\title{Appunti di Lab II}
\author{Tommaso Miliani}
\date{11-03-26}

\begin{document}
\newtheoremstyle{theoremEnv}
                {}          % Space above
                {}          % Space below
                {\slshape}  % Body font
                {}          % Indent amount
                {\bfseries} % Head font
                {.}         % Punctuation after head
                {\newline}  % Space after theorem head
                {}          % Theorem head spec
\theoremstyle{theoremEnv}

\newtheorem{definition}{Definizione}[section]
\newtheorem{theorem}{Teorema}[section]
\newtheorem{lemma}{Proposizione}[section]
\newtheorem{observation}{Osservazione}[section]
\newtheorem{corollary}{Corollario}[theorem]
\newtheorem{example}{Esempio}[section]
\newtheorem{remark}{Enunciato}[section]

\maketitle

\section{Punchlines}
\begin{itemize}
    \item Porca Miseria
    \item Siamo del Gatto
    \item A questo punto possiamo mangiare la foglia
    \item Porca miseria!
    \item Poverino questo generatore
    \item Questo *BAM BAM* multimetro misura tutto
    \item La molla dopo un po' smolla
    \item Quello del pi greco era grullo
    \item Avete mangiato la foglia
\end{itemize}

\section{Misurazioni di corrente: il galvanometro}
Il Galvanometro è lo strumento che permette di eseguire 
misure di corrente e tensione. Per poter misurare la corrente,
si prende una spira di area $S$ ed un circuito: sottoponendo la spira ad un campo 
magnetico esterno, in cui circola della corrente $i$, può essere descritta dal 
vettore momento magnetico come 
\begin{gather*}
    \vv{m} = iS\hat{n}  
\end{gather*} 
Dove $\hat{n}$ è il versore normale uscente dal circuito. Se si ponesse un 
campo magnetico sulla spira, essa sentirà un momento torcente di
\begin{gather*}
    \vv{M} = \vv{m} \times \vv{B}   
\end{gather*} 
La spira inizia dunque a roteare proporzionalemente alla corrente stessa. 
Se ci si mettesse una molla tale per cui si possa avere un momento meccanico 
sullo stesso asse lungo il momento magnetico
\begin{gather*}
    \vv{M_M} = -k\phi \hat{k}  
\end{gather*}
Si ha dunque uno strumento che dà una risposta proporzionale all'aumento di corrente.
\begin{wrapfigure}{r}{0.4\textwidth}
    \centering
    \caption{}
    \begin{circuitikz}
        \draw(0, 0) to [short, *-] (0, -1)
        to [R, l_=$\rho_G$] (0, -2) 
        to [short] (0, -2.5)
        to [rmeter, l_=$G$] (0, -3.5)
        to [short, -*] (0, -4);
    \end{circuitikz} 
\end{wrapfigure}
Allora si può misurare la corrente come 
\begin{gather*}
    i = k_R \phi 
\end{gather*}
Dove $k_R$ è una costante che determina la precisione del galvanometro che può arrivare
ad una precisione di $10^{-6}$ Ampere su radianti. Con un sistema di leva ottica, alla fine la 
sensibilità di questo oggetto può anche arrivare a $100$ nanoAmpere per divisione, dove la divisione è
la divisione della scala graduata. 
Generalmente questo strumento ha una resistenza molto bassa e permette dunque
di poter misurare tutto in un circuito. Attualmente i valori di riferimento per cui si 
definisce la corrente sono definiti su standard quantistici. La curva di taratura ideale 
dello strumento si determina con l'incertezza sia sulla lettura che sul
valore effettivamente misurato, dove $\frac{1}{k_R}$ è il coefficiente della
curva di taratura (retta):
\begin{gather*}
    \begin{tikzpicture}
        \draw[->](0, 0) -- (3, 0) node[at end, below] {$i$};
        \draw[->](0, 0) -- (0, 3) node[at end, left] {$\phi$};
        \draw(0, 0) -- (3, 3) node[at end, right] {$\frac{1}{k_R}$};
        \draw[|-|](0.5, 0.7) -- (0.5, 0.3);
        \draw[|-|](0.3, 0.5) -- (0.7, 0.5);
        \draw[|-|](1.5, 1.7) -- (1.5, 1.3);
        \draw[|-|](1.3, 1.5) -- (1.7, 1.5);
        \draw[|-|](2.5, 2.7) -- (2.5, 2.3);
        \draw[|-|](2.3, 2.5) -- (2.7, 2.5);
    \end{tikzpicture}
\end{gather*}
Anche negli strumenti digitali il concetto è lo stesso poiché l'errore
di lettura può essere dato dal numero di bit del sistema digitale
(sull'ultima cifra, o anche le ultime due) le cifre ballano. \\
È ovvio che il Galvanometro abbia una resistenza interna molto piccola. Se infatti si prendesse
un circuito qualunque con un generatore ed una resistenza e ci si attaccasse un galvanometro,
si dovrebbe considerare le due correnti
\begin{gather*}
    i = \frac{\epsilon}{\rho} \qquad i_m = \frac{\epsilon}{\rho + \rho_G}
\end{gather*}
Che sono ovviamente diverse, saranno simili se e solo se $\rho_G << \rho$.

\begin{wrapfigure}{r}{0.4\textwidth}
    \centering
    \caption{}
    \begin{circuitikz}
        \draw(0, 0) to [short, *-] (1, 0);
        \draw(1, 0) to [short, ] (1, 1)
        to [R, l_=$\rho_G$] (2.5, 1) 
        to [rmeter, l_=$\rho_G$] (4, 1)
        to [short] (4, 1) to [short] (4, 0);
        \draw(1, 0) to [short] (1, -1)
        to [R, l_= $\frac{\rho_G}{9}$] (4, -1)
        to [short] (4, -1) 
        to[short] (4, 0);
        \draw(4, 0) to [short, -*] (5, 0);
    \end{circuitikz}    
\end{wrapfigure}
Nei casi in cui $\rho_G$ non sia sufficientemente piccola (o nei casi in cui la corrente è 
talmente grande che si arriverebbe a fondo scala) si utilizza il teorema dell'"asciugacapelli 
in bagno coi piedi bagnati": la corrente passa sempre dove vuole e dove c'è meno corrente.
Sostanzialmente si utilizza lo \textbf{shunt} di corrente: in parallelo al galvanometro 
si prende una resistenza più piccola di $\rho_G$. 
\begin{gather*}
    \rho = \frac{\rho_G}{9} \qquad I = i_S + i_G
\end{gather*}
Allora si avrà che la corrente che passa nello shunt e nel galvanometro è
\begin{gather*}
    \frac{i_G}{i_S} = \frac{\rho_G}{9\rho_G} = \frac{1}{9}
\end{gather*}
Allora la sommad delle due correnti:
\begin{gather*}
    9i_G + i_G = I = 10 i_G \ \Longrightarrow \ \boxed{i_G = \frac{1}{10}I}
\end{gather*}
In generale lo shunt è come affettare un blocco, perdendo però sensibilità. Ponendo questo circuito 
nel circuito del quale si vuole misurare la corrente, la resistenza equivalente del circuito è
\begin{gather*}
    R_T = \frac{\rho_G \frac{\rho_G}{9}}{\rho_G + \frac{\rho_G}{9}}  =\frac{\rho_G}{10} 
\end{gather*}
Questo va bene perché qualsiasi effetto sistematico si sia introdotto con il galvanometro è ridotto di un 
fattore 10:
\begin{gather*}
    i_m = \frac{\epsilon}{\rho + \frac{\rho}{10}} \sim i
\end{gather*}

\section{Divisore di tensione}
\begin{wrapfigure}{r}{0.4\textwidth}
    \centering
    \caption{Potenziometro}
    \begin{tikzpicture}
        
    \end{tikzpicture}    
\end{wrapfigure}
Il divisore di tensione è un oggetto che divide la tensione in ingresso 
con il rapporto impostato dalle manopole su di esso. Un modo per realizzare un 
divisore di tensione è quello di realizzare un \textbf{potenziometro}: esso è un contatto strisciante
che permette di modulare la fem in ingresso secondo la seguente relazione
\begin{gather*}
    \frac{l\rho}{S} = R
\end{gather*}
Per avere un fattore di un millesimo, $R(l(\phi))$ dovrebbe essere $\frac{1}{1000}R$, 
altrimenti si può prendere 1000 resistenze tutte uguale, tuttavia ci deve essere un modo 
più furbo per poter realizzare un divisore di tensione.
La linearità è uguale a
\begin{gather*}
    x \% \text{tensione in ingresso}
\end{gather*}
Ogni manopola ha una divisione che è uguale alla tensione 
in ingresso per quella in uscita:
\begin{gather*}
    v_{out} = rV
\end{gather*}
Allora l'errore di linearità è
\begin{gather*}
    \frac{\Delta v}{v} = \frac{\text{linearità} \cdot V}{v} = \frac{\text{linearità} \%}{r}
\end{gather*}
In laboratorio si ha una linearità di $20$ parti per milione ed un errore di linearità pari a
\begin{gather*}
    \frac{\Delta v}{v} = 40 \ ppm \approx 4 \cdot 10^{-6}
\end{gather*} 

\subsection{Realizzare un divisione di tensione}
Il divisore di tensione Kelvin-Varley è un divisore di tensione che presenta 
delle tensioni in ingresso ed in uscita con rapporto di divisone $r$ con divisone 
3 decadi (ossia 3 sole manopole che si possono girare). 
\begin{gather*}
    \begin{circuitikz}
        \draw(0, 0) to [short] (1, 0);
    \end{circuitikz}
\end{gather*}

\begin{wrapfigure}{r}{0.4\textwidth}
    \centering
    \caption{Divisore di tensione}
    \begin{circuitikz}
        
    \end{circuitikz}    
\end{wrapfigure} \noindent
A partire da $H$, si utilizzano un certo numero di resistenze e si arriva ad $L$ 
con una serie di resistenze (11 da 1Kilo $\Omega$); a partire dalla prima resistenza (che non si conta perché 
è molto piccola e serve solo a contrastare la resistenza del conduttore). Queste 11 resistenze tutte 
uguali sono messe in modo tale che la manopola 1 possa prendere sia la posizione $i$ che la posizione $i + 2$ esima.
Dopo la posizione $i$ e $i + 2$ si pone un altro circuito con 11 resistenze
scelte in modo tale che siano un quinto di quelle precedenti (200 $\Omega$), si prende di nuovo una manopola
che seleziona $j$ e $j + 2$ e fa passare la corrente lungo un circuito con 10 resistenze
che abbiano un quinto della resistenza della decade precedente (40 $\Omega$). L'ultima decade 
ha un solo cursore e seleziona una resistenza $k$ in modo tale che la corrente arrivi ad $h$. Dalla resistenza
zero prima della prima decade, essa arriva ad $l$. \\
La resistenza in ingresso è la stessa in uscita a prescindere da tutti i cursori infatti,
indipendentemente dal selettore, il divisore è fatto in modo che si abbiano sempre
in parallelo due resistenze uguali, dunque la resistenza che è in parallelo alla 
prima decade è $200\Omega$ per 10, la quale è esattamente la resistenza 
compresa tra i due selettori della priam decade. Allora la resistenza
tra $h$ e $l$ è data da 2$k\Omega$ in parallelo a $2k\Omega$ è $1k\Omega$, 
la quale è in serie con le altre 9 resistenze, dunque la resistenza totale è sempre $10 k\Omega$
indipendentemente dal selettore.  \\
Si vuole adesso trovare la differenza tra uscita ed entrata:
\begin{gather*}
    v_h - v_l = f(V_H - V_L) \ \Longrightarrow \ v_h - v_l = v_h - v_F + v_F - v_D + v_D - v_l
\end{gather*}
In modo da valutare le tre differenze di potenziale a destra. 
\begin{gather*}
    v_D - v_l = \frac{i \cdot 1k\Omega}{10 k\Omega} (V_H - V_L) = \frac{i}{10} (V_H - V_L)
\end{gather*}
Dove $i$ è l'indice di selettore della prima decade. 
\begin{gather*}
    v_F - v_D = \frac{j}{100}(V_H - V_L) \qquad v_h - v_F = \frac{k}{1000} (V_H - V_L)
\end{gather*}
Allora se si impostasse una decade su 5, l'altra su 3, e l'ultima su 2, la tensione in uscita è
moltiplicata per $0.532$. Si dimostrano ora le altre due. Per la seconda, indipendentemente 
dal selettore la resistenza totale tra $F$ e $V$ sarà sempre $10 \cdot 200\Omega$, allora,
rispetto alla tensione in ingresso
\begin{gather*}
    v_F - v_D = \frac{j \cdot 200\Omega}{10 \cdot 200 \Omega} (V_C - V_D)
\end{gather*}
Dove $(V_C - V_D)$ si lega a $(V_H - V_L)$ nella seguente maniera:
\begin{gather*}
    V_C - V_D = \frac{1}{10}(V_H - V_L) \ \Longrightarrow \ v_F - v_D = \frac{j}{100}(V_H - V_L) 
\end{gather*}
Lo stesso conto si fa per l'ultimo.


\end{document}